The game that we modified makes the connection between inference and planning measurable through
explicit information sets and verified action values. We want to discuss three thought.

\paragraph{Richer partner states.}
Our diagnosis concentrates uncertainty on a partner's preferences. In less
constrained interactions, an agent must also infer what its partner has
observed, what it believes others know, and which
behavioral rule it follows. The same refusal, for example, could express an
unfavorable preference or uncertainty about information unavailable to the
partner. Recent work on multi-agent theory of mind examines private knowledge,
adaptive reasoning about another agent's mental state, and hypotheses about
partner strategies \citep{ys2026sotopiatom,mu2026adaptive,cross2025hypothetical}.
These are extensions of \emph{inference} as defined here: each concerns a
latent state needed to predict a partner's response and plan one's own action.
Extending game slices to such states would test whether the separation between
inference and planning remains useful when beliefs themselves have structure.

\paragraph{Strategic coverage.}
The appeal of game slices lies in control over which decisions become training
examples. Complete-game self-play generates coherent trajectories with the
learner's current partners, but its training distribution is the set of nodes
those policies actually reach. A rarely visited offer, rejection, or
information-gathering decision may therefore receive little training signal
even when it is strategically decisive. We instead construct and sample
slices around specified interaction patterns, retaining cases for which the
reference computation supplies verified inference and action-value targets.
Work on subgame curricula likewise shows that changing the distribution of
starting states can accelerate multi-agent learning \citep{chen2024subgame}.
This offers a possible explanation for the efficiency of targeted slices,
though our comparison does not isolate state selection from denser
supervision. The complementary strengths suggest a concrete hybrid: use
complete games to discover consequential failures, solve and train on slices
around those decisions, then return to full interactions to test whether the
learned behavior survives its wider context.

\paragraph{Robustness across partner policies.}
Our reference labels assume one selected partner policy, while self-play
primarily exposes the learner to copies of its evolving policy. This matters
for inference as well as action selection: an offer that is impossible under
one reference policy may be ordinary for a partner with a different planning
horizon, tie rule, or degree of noise. Evidence that self-play agents can
learn conventions that fail with unfamiliar partners, and that diverse
training partners can improve cross-partner coordination motivates testing
policy variation directly
\citep{hu2020otherplay,strouse2021collaborating}. One direction is to train
against a population of partner policies and generate slice labels under
each policy, including controlled deviations from the reference solver.
Belief targets could then marginalize over uncertainty about both preferences
and behavioral policies, while evaluation holds out entire policy families.
Such tests would distinguish strategic behaviors that transfer across
partners from responses tailored to the particular oracle used here.
